% *==================================================================================*
% *                     Review vs. Camera-Ready settings                             *
% *==================================================================================*
%
% REVIEW: Use the following command for submitting the paper (double-blind,
% for review):
% \documentclass{Interspeech}
%
% CAMERA-READY: Use the following command for the camera-ready version, one
% affiliation per line:
\documentclass[cameraready]{Interspeech}
% *==================================================================================*


% **************************************
% *                                    *
% *      STOP !   DO NOT DELETE !      *
% *          READ THIS FIRST           *
% *                                    *
% * This template also includes        *
% * important INSTRUCTIONS that you    *
% * must follow when preparing your    *
% * paper. Read it BEFORE replacing    *
% * the content with your own work.    *
% **************************************

%==================================================================================
% Title
% Must exactly match the title entered into the paper submission system
\title{Distilling Structured Reasoning into SpeechLLMs for Spoken Language Understanding}

%==================================================================================
% Authors
% The order of authors here must exactly match the order entered into the paper submission system
% Note that the COMPLETE list of authors MUST be entered into the paper submission system at the outset, including when submitting your manuscript for double-blind review
% The ORCID number is still optional but will become mandatory in the future years. It is strongly encouraged to get an ORCID for each cu-author.
% Middle names, including initials, must be included in the first name
\author[affiliation={1}]{Tsukagoshi}{Toshihiro}
\author[affiliation={2}]{Natsuo}{Yamashita}
\author[affiliation={1,3}]{FirstNameC}{LastNameC}
% The maximum number of authors in the author list is 20. If the number of contributing authors is more than this, they should be listed in a footnote or the acknowledgement section.

%==================================================================================
% Affiliations

\address{
    $^1$ Shizuoka University, Japan \\
    $^2$ Address Affiliation 2, Country Affiliation 2 
}

%==================================================================================
% Emails
\email{tsukagoshi.toshihiro.20@shizuoka.ac.jp, second@companyA.com, third@companyB.ai}

%==================================================================================
% Keywords
\keywords{spoken language understanding, speech large language models, chain-of-thought reasoning, multitask learning}

\newcommand{\blue}[1]{\textcolor{blue}{#1}}
\newcommand{\red}[1]{\textcolor{red}{#1}}
\usepackage{comment}

%==================================================================================
% Content

\begin{document}

\maketitle

% the abstract here must exactly match the abstract entered into the paper submission system
\begin{abstract}
Spoken language understanding (SLU) maps speech signals to structured semantic representations, such as intent classification and slot filling. Recent Speech Large Language Models (SpeechLLMs) enable end-to-end SLU by directly processing acoustic inputs. However, standard fine-tuning often exhibits confusion among semantically adjacent intent categories, resulting in insufficient separation between competing interpretations. To mitigate this issue, we propose a structured reasoning distillation framework for SpeechLLM-based SLU. We train the model to learn a stepwise decision process that explicitly contrasts plausible candidate interpretations, rejects incorrect alternatives, and produces the final structured prediction. We generate supervision for this intermediate decision process using a large reasoning model and distill it into SpeechLLMs. Since smaller models may produce unstable structured outputs, we adopt a multitask learning formulation that jointly optimizes intermediate reasoning and final prediction. Experiments on SLURP across six architectures demonstrate consistent improvements over standard fine-tuning.

\end{abstract}






\section{Introduction}

Spoken language understanding (SLU) is a key component of task-oriented dialogue systems, enabling a wide range of real-world applications such as commercial voice assistants~\cite{SaxoChoudharyMcKennaMouchtaris2021E2ESLUVoiceAssistants}, customer-service automation in contact centers~\cite{WuMaoFeng2022AIOnlineCustomerService}, and supporting front-line workers in operational environments~\cite{LinaresGarcia2022VIVA}. 
SLU typically comprises intent classification~\cite{Atuhurra2024DomainAI} and slot filling~\cite{weld2022survey} for mapping an utterance to a semantic frame.


Recent advances in Speech Large Language Models (SpeechLLMs), including SALMONN~\cite{tang2024salmonn}, Qwen2-Audio~\cite{chu2024qwen2audio}, Qwen2.5-Omni~\cite{xu2025qwen25omni}, and Audio-Flamingo~\cite{goel2025audioflamingo3}, enable end-to-end SLU by directly processing acoustic signals and generating structured textual outputs.


Existing studies leveraging SpeechLLMs for SLU have primarily focused on adaptation and generalization.
Some works explore zero-shot or prompting-based SLU by reformulating intent detection and slot filling as question-answering tasks~\cite{LiKeizerDoddipatla2024PromptingWhisperQASLU}, while others evaluate large-scale zero-shot slot filling across diverse schema settings~\cite{hacioglu2025speechllms,hacioglu2025slot}.
Supervised adaptation approaches include parameter-efficient fine-tuning with modality alignment~\cite{li2024whisma}, task-agnostic generalization to unseen intents~\cite{AgrawalGanapathy2025SLUUnseenICL}, and systematic analyses of fine-tuning strategies under limited speech supervision~\cite{choi2025exploring}.
While prior studies have explored the potential of SpeechLLMs for SLU, errors still frequently occur between semantically adjacent intents that differ only in subtle aspects such as scope or granularity.
Learning such fine-grained distinctions purely from data typically requires a substantial number of boundary examples. 
However, collecting sufficiently diverse and finely differentiated training data is costly in practice. Therefore, methods that can efficiently learn decision boundaries from limited data are needed.


An effective alternative to data-intensive boundary learning is to explicitly supervise the intermediate decision process.
In text-based LLM research, rationale-based supervision~\cite{zelikman2022starbootstrappingreasoningreasoning} and distillation methods~\cite{hsieh2023distillingstepbystepoutperforminglarger} have demonstrated that incorporating structured intermediate reasoning can improve downstream performance.
Recent reasoning distillation studies~\cite{magister-etal-2023-teaching, HoSchmidYun2023ReasoningTeachers, Fu2022ComplexityBasedPF} further show that transferring reasoning traces from larger teacher models to smaller students can yield measurable gains under limited model capacity.

However, subsequent analyses suggest that the effectiveness of reasoning supervision depends critically on how it is incorporated.
Naively injecting detailed reasoning traces into smaller models does not always lead to improvements, and excessive reasoning granularity may even hinder optimization~\cite{li-etal-2025-small-models, chen-etal-2025-unveiling-key, chen-etal-2025-skip, hsieh-etal-2023-distilling}.
These findings indicate that reasoning supervision should be treated not as a standalone objective, but as a carefully designed auxiliary signal that shapes the representation space.

Moreover, prior studies have focused primarily on text-only settings, leaving open the question of how structured reasoning supervision behaves in SpeechLLM-based SLU, where acoustic and semantic representations must be jointly modeled.
Motivated by these considerations, we incorporate rationale distillation as an auxiliary objective within a multitask learning framework, allowing structured reasoning to regularize intent prediction without destabilizing generation.



In this work, we investigate whether structured reasoning supervision can improve intent discrimination in moderate-scale SpeechLLMs for end-to-end SLU.
Rather than relying on emergent chain-of-thought behavior, we explicitly model SLU as a stepwise decision process that contrasts plausible intent candidates, rejects incorrect alternatives with evidence-grounded justification, and finally produces a structured semantic label.
We generate such structured rationales using a large reasoning teacher model and distill them into SpeechLLMs under a multitask learning framework, where rationale generation serves as an auxiliary regularizer for direct intent prediction.
Through comprehensive experiments on SLURP across six architectures, we demonstrate that this structured reasoning distillation consistently improves intent accuracy and SLU-F1, while also enlarging inter-class margins in the latent representation space.





\section{Proposed Method}
\label{sec:method}

The overall framework of our approach is illustrated in 
Figure~\ref{fig:framework}. We augment moderate-scale SpeechLLMs 
with structured reasoning to improve performance in low-resource 
SLU settings.

\subsection{Tripartite Rationale Design: $C \to R \to L$}

To transform implicit SLU reasoning into an explicit supervisory 
signal, we define a structured rationale sequence 
$T_{cot} = \{C, R, L\}$. Each component corresponds to a distinct 
reasoning step that mirrors human decision-making under label 
ambiguity.

\begin{itemize}
    \item \textbf{Intent Candidates ($C$)}: Models fine-tuned for SLU 
    often confuse semantically adjacent intents (e.g., 
    \textit{calendar\_set} vs.\ \textit{reminder\_set}), as a 
    direct classification objective provides no explicit incentive 
    to contrast them. To sharpen decision boundaries, we first 
    present a shortlist of plausible intent labels filtered from 
    domain metadata, encouraging contrastive reasoning among the 
    most confusable classes before committing to a prediction.

    \item \textbf{Rejection Rationale ($R$)}: Enumerating candidates 
    alone is insufficient if the model cannot articulate why one is 
    preferred over others. We therefore require a natural language 
    explanation grounded in acoustic or semantic evidence, 
    explicitly justifying the rejection of competing candidates. 
    This step encourages attention to discriminative cues rather 
    than reliance on surface-level lexical patterns.

    \item \textbf{Final Structured Label ($L$)}: After resolving the 
    ambiguity through explicit contrast and justification, the model 
    outputs the final structured semantic representation corresponding 
    to the predicted intent and slot values.
\end{itemize}


By structuring the reasoning in this order, the model is guided 
to progressively narrow the hypothesis space through explicit 
contrastive reasoning, rather than collapsing the decision into 
a single opaque mapping — ultimately producing intent 
representations with greater inter-class separability.
\subsection{Multitask CoT Framework}

A key limitation of moderate-scale SpeechLLMs is that structured 
reasoning does not emerge naturally from scale alone. Rather than 
waiting for such emergence, we explicitly inject reasoning ability 
as a supervised objective through a multitask formulation.

Concretely, the model is conditioned via a mode tag prepended to 
the input, directing it to operate in one of two modes:

\begin{itemize}
    \item \textbf{CoT-mode}: The model generates the full rationale 
    sequence $T_{cot} = \{C, R, L\}$, explicitly performing candidate 
    contrast and rejection before producing the final structured label.

    \item \textbf{Label-mode}: The model directly outputs the final 
    structured label $T_{label} = \{L\}$, without generating intermediate 
    rationale components.
\end{itemize}



Both modes share identical model parameters; the distinction lies 
solely in the generation target specified by the mode tag. Through 
joint training, the reasoning patterns internalized via CoT-mode 
act as a regularizer that shapes the shared representation space, 
enabling Label-mode to produce more discriminative predictions 
without relying on long-form rationale generation at inference time.


\section{Experimental Setup}
\label{sec:experiments}

\subsection{Dataset}
We evaluate our framework on the SLURP 
dataset~\cite{bastianelli-etal-2020-slurp}, which comprises 
11,514 training, 2,033 development, and 2,974 test utterances 
covering 18 domains and 60 intent classes. Following 
\cite{choi2025exploring}, all models are trained on both speech 
and gold transcription inputs to facilitate acquisition of 
fundamental SLU capabilities. For Direct and Standalone CoT 
models, training proceeds for 2 epochs. For Multitask CoT, 
since the model is jointly supervised on both CoT-mode and 
Label-mode targets, the effective number of parameter updates 
per sample is doubled; we therefore train for 1 epoch to ensure 
a fair comparison in terms of total training steps across all 
conditions.

\subsection{Metrics}
Performance is assessed across four complementary metrics. 
Scenario and Action accuracy measure the correctness of 
high-level domain classification (e.g., \textit{calendar}, 
\textit{email}) and specific intent action (e.g., \textit{set}, 
\textit{query}), respectively. Intent accuracy evaluates the 
combined (Scenario, Action) pair, reflecting overall 
classification correctness. SLU-F1~\cite{bastianelli-etal-2020-slurp} 
provides a holistic measure of the full semantic frame, 
penalizing slot errors more heavily when the intent is incorrect.


\begin{enumerate}
    \item \textbf{Scenario}: Accuracy of the high-level domain classification (e.g., \textit{calendar}, \textit{email}).
    \item \textbf{Action}: Accuracy of the specific intent action (e.g., \textit{set}, \textit{query}).
    \item \textbf{Intent}: Accuracy of the combined (Scenario, Action) pair. This represents the overall classification correctness.
    \item \textbf{SLU-F1}: A comprehensive metric that evaluates the entire semantic frame. Unlike standard F1, SLU-F1~\cite{bastianelli-etal-2020-slurp} penalizes slot errors more heavily if the intent is incorrect, providing a holistic measure of the model's understanding of both global intent and local entities.
\end{enumerate}
\begin{table*}[t]
\centering
\caption{Main results on the SLURP test set. Each column group reports 
Scenario (Sce), Action (Act), Intent (Int), and SLU-F1 (F1) accuracy. 
The proposed Multitask CoT consistently achieves the best performance 
across all architectures.}
\label{tab:main_results}
\resizebox{\textwidth}{!}{
\begin{tabular}{lcccccccccccc}
\hline
 & \multicolumn{4}{c}{Direct} 
 & \multicolumn{4}{c}{Standalone CoT (CRL)} 
 & \multicolumn{4}{c}{Multitask CoT (Proposed)} \\ 
\cmidrule(lr){2-5}\cmidrule(lr){6-9}\cmidrule(lr){10-13}
Model & Sce & Act & Int & F1 
      & Sce & Act & Int & F1 
      & Sce & Act & Int & F1 \\ \hline
Qwen2.5-Omni-7B~\cite{xu2025qwen25omni}   & 91.53 & 88.74 & 87.83 & 76.20 & 91.12 & 87.90 & 87.05 & 76.08 & \textbf{91.96} & \textbf{88.97} & \textbf{88.23} & \textbf{77.10} \\
Qwen2.5-Omni-3B~\cite{xu2025qwen25omni}  & 90.65 & 87.79 & 86.79 & 75.04 & 89.27 & 85.47 & 84.47 & 73.23 & \textbf{91.39} & \textbf{88.77} & \textbf{87.79} & \textbf{76.30} \\
Qwen2-Audio-7B~\cite{chu2024qwen2audio}     & 89.37 & 85.21 & 84.23 & 70.08 & 88.63 & 84.80 & 83.46 & 70.62 & \textbf{90.15} & \textbf{86.65} & \textbf{85.71} & \textbf{72.22} \\
Qwen2-Audio-7B-Instruct~\cite{chu2024qwen2audio}   & 88.67 & 85.84 & 84.70 & 70.83 & 88.84 & 85.00 & 83.73 & 71.42 & \textbf{90.11} & \textbf{87.19} & \textbf{86.38} & \textbf{72.79} \\
Audio-Flamingo-3~\cite{goel2025audioflamingo3}    & 87.93 & 84.10 & 82.65 & 69.98 & 87.90 & 84.33 & 83.05 & 71.92 & \textbf{89.88} & \textbf{86.45} & \textbf{85.44} & \textbf{72.86} \\
Music-Flamingo~\cite{ghosh2025musicflamingoscalingmusic}   & 87.69 & 83.89 & 82.65 & 69.02 & 87.96 & 84.36 & 82.92 & 71.42 & \textbf{90.01} & \textbf{86.65} & \textbf{85.74} & \textbf{72.11} \\ \hline
\end{tabular}
}
\end{table*}

\subsection{Representation Analysis}

To quantitatively assess whether multitask CoT supervision improves 
intent discriminability, we extract hidden state vectors at the moment 
of intent generation. Specifically, we identify the tokens corresponding 
to the scenario and action values in the generated output and pool their 
hidden states to obtain a single intent representation vector per 
utterance:

\begin{equation}
    \mathbf{v} = \frac{1}{|\mathcal{T}|} \sum_{t \in \mathcal{T}} \mathbf{h}_t
\end{equation}

\noindent where $\mathcal{T}$ denotes the set of token positions 
corresponding to the scenario and action values, and $\mathbf{h}_t$ 
is the hidden state at position $t$ from the final layer.

We evaluate inter-class separability using two complementary metrics. 
The \textbf{Silhouette coefficient} for sample $i$ is defined as:

\begin{equation}
    s(i) = \frac{b(i) - a(i)}{\max\{a(i),\, b(i)\}}
\end{equation}

\noindent where $a(i)$ is the mean cosine distance to samples in the 
same intent class, and $b(i)$ is the mean cosine distance to samples 
in the nearest competing class. The overall score $\bar{s} = 
\frac{1}{N}\sum_i s(i)$ approaches $1$ when clusters are 
well-separated.

The \textbf{nearest-centroid margin} for sample $i$ is defined as:

\begin{equation}
    m(i) = \min_{c \neq y_i} \lVert\mathbf{v}_i - \boldsymbol{\mu}_c\rVert_2 
           - \lVert\mathbf{v}_i - \boldsymbol{\mu}_{y_i}\rVert_2
\end{equation}

\noindent where $\boldsymbol{\mu}_c$ is the centroid of intent class 
$c$ and $y_i$ is the ground-truth intent of sample $i$. A larger 
positive margin indicates a clearer decision boundary between 
semantically adjacent intents.



\subsection{Implementation Details}
All models are fine-tuned using BF16 precision with the AdamW 
optimizer, a learning rate of $4 \times 10^{-5}$, a batch size 
of 4, and 4 gradient accumulation steps. Oracle rationales for 
multitask training are generated via DeepSeek-R1, whose 
chain-of-thought pretraining makes it well-suited for producing 
structured, evidence-grounded reasoning traces in the 
$C \to R \to L$ format. For Direct and Standalone CoT models, 
training proceeds for 2 epochs. For Multitask CoT, since the 
model is jointly supervised on both CoT-mode and Label-mode 
targets, the effective number of parameter updates per sample 
is doubled; we therefore train for 1 epoch to ensure a fair 
comparison in terms of total training steps across all conditions.

\begin{table}[t]
\centering
\caption{Representation geometry metrics for Qwen2.5-Omni-7B. 
Higher values indicate greater inter-class separability.}
\label{tab:rep_analysis}
\resizebox{\columnwidth}{!}{
\begin{tabular}{lccc}
\hline
\textbf{Metric} & \textbf{Direct} & \textbf{Multitask CoT} & \textbf{Diff} \\ \hline
Silhouette (cosine)          & 0.242 & 0.279 & +0.037 \\
Silhouette (euclidean)       & 0.157 & 0.176 & +0.019 \\
Fisher ratio                 & 1.133 & 1.276 & +0.143 \\
Inter-centroid cosine dist.  & 0.262 & 0.328 & +0.067 \\
Norm. centroid margin (mean) & 0.174 & 0.194 & +0.020 \\ \hline
\end{tabular}
}
\end{table}

\section{Results}
\label{sec:results}

\subsection{Main Results}
Table~\ref{tab:main_results} reports Scenario, Action, Intent accuracy and SLU-F1 for three training conditions across six architectures: Direct fine-tuning, Standalone CoT, and the proposed Multitask CoT using \texttt{<Label>} mode inference.
Multitask CoT consistently outperforms both baselines across all models and metrics, with intent accuracy gains ranging from +0.96\% (Qwen2.5-Omni-7B) to +2.79\% (Qwen2.5-Omni-3B) over Direct fine-tuning, and SLU-F1 improvements of up to +1.26 points (Qwen2.5-Omni-7B).
Notably, Standalone CoT occasionally underperforms Direct fine-tuning, with intent accuracy degrading by up to 2.32\% (Qwen2.5-Omni-3B), indicating that CoT supervision without a complementary \texttt{<Label>} mode objective does not reliably improve performance in moderate-scale models.
These results demonstrate that explicitly supervising reasoning as a multitask objective is more effective than relying on emergent CoT behavior for moderate-scale SpeechLLMs.


\subsection{Representation Analysis}

Table~\ref{tab:rep_analysis} reports representation geometry 
metrics comparing Direct fine-tuning and Multitask CoT 
(Qwen2.5-Omni-7B). Multitask CoT improves across all 
separability metrics: the Silhouette coefficient (cosine) 
increases from 0.242 to 0.279 (+0.037), and the Fisher ratio 
rises from 1.133 to 1.276 (+0.143), indicating greater 
between-class scatter relative to within-class scatter. The 
mean inter-centroid cosine distance also increases from 0.262 
to 0.328 (+0.067), suggesting that intent clusters are more 
evenly spread across the representation space. Furthermore, 
the normalized nearest-centroid margin improves from 0.174 to 
0.194 (+0.020), confirming that decision boundaries between 
semantically adjacent intents become more clearly defined 
under multitask CoT supervision.



\subsection{Impact of Rationale Components}
Table~\ref{tab:ablation_components} reports intent accuracy and SLU-F1 for each rationale component configuration using \texttt{<Label>} mode inference. Among individual components, C yields the largest improvement in intent accuracy (+0.37\% over baseline), indicating that explicitly enumerating candidate intents contributes most to sharpening decision boundaries. R alone provides marginal gains, suggesting that rejection rationale is insufficient as a standalone supervisory signal. The full configuration C+R+L achieves the highest intent accuracy, while C+L and C+R+L both attain the highest SLU-F1 of 77.18, indicating that the combination of candidate enumeration and rejection rationale contributes to stable performance across metrics.



\begin{table}[t]
\centering
\caption{Ablation of rationale components (Qwen2.5-Omni-7B, Label-branch inference). \textbf{C}: Intent Candidates, \textbf{R}: Rejection Rationale, \textbf{L}: Label (direct prediction).}
\label{tab:ablation_components}
\begin{tabular}{lcccc}
\hline
\textbf{Configuration} & \textbf{Sce / Act / Int} & \textbf{SLU-F1} \\ \hline
Direct (Baseline)      & 91.53 / 88.74 / 87.83 & 76.20 \\
Multitask w/ C         & 91.56 / 89.00 / 88.20 & 77.18 \\
Multitask w/ R         & 91.53 / 88.84 / 87.96 & 76.01 \\
Multitask w/ C+R       & 91.83 / 88.87 / 88.13 & \textbf{77.18} \\
\textbf{Multitask w/ C+R+L} & \textbf{91.96 / 88.97 / 88.23} & 77.10 \\ \hline
\end{tabular}
\end{table}

The results indicate that providing Intent Candidates ($C$) yields the most significant initial boost in intent accuracy (+0.37\% over baseline). This supports our hypothesis that forcing the model to explicitly contrast similar intent categories early in the process sharpens the decision boundaries. While Rejection Rationale ($R$) alone provides marginal gains, its combination with candidates (CRL) achieves the highest overall accuracy and stability across domains.

\subsection{Inference Branch: CoT vs. Label}
Table~\ref{tab:branch_comparison} reports intent accuracy and SLU-F1 for each multitask configuration under two inference modes: \texttt{<CoT>} and \texttt{<Label>}. Across all configurations, \texttt{<Label>} mode consistently achieves higher performance than \texttt{<CoT>} mode. Notably, the performance gap is particularly pronounced in the Multitask w/ R+L configuration, where \texttt{<Label>} mode outperforms \texttt{<CoT>} mode by 2.06\% in intent accuracy. This trend suggests that CoT-mode supervision effectively shapes the shared representation space during training, but the benefit does not transfer to inference-time rationale generation. Rather, the structured reasoning internalized through CoT-mode training is more effectively utilized when the model performs direct prediction via \texttt{<Label>} mode, without the overhead of generating intermediate rationale sequences.


\section{Conclusion}
\label{sec:conclusion}
In this paper, we proposed a Multitask CoT framework for 
moderate-scale SpeechLLMs to address the challenges of spoken 
language understanding. By supervising a structured reasoning process 
($C \to R \to L$) alongside direct intent prediction, we successfully 
improved performance across six distinct model architectures. 
Our representation analysis, including Silhouette coefficients and 
Fisher ratio, provided quantitative evidence that multitask rationale 
supervision encourages the model to develop more separable intent 
clusters. Crucially, we demonstrated that using the rationale as a 
training regularizer while performing direct inference (\texttt{<Label>} 
mode) allows the model to leverage the benefits of reasoning without 
the overhead or noise of long-form generation. Future work will explore 
the scalability of this approach to even smaller edge-device models and 
its application to multilingual SLU scenarios.

\begin{table}[t]
\centering
\caption{Comparison of inference modes within the multitask framework.}
\label{tab:branch_comparison}
\begin{tabular}{llcc}
\hline
\textbf{Configuration} & \textbf{Mode} & \textbf{Intent Acc.} & \textbf{SLU-F1} \\ \hline
Multitask w/ C+L   & \texttt{<CoT>}   & 87.66 & 77.31 \\
                   & \texttt{<Label>} & \textbf{88.06} & \textbf{77.89} \\ \hline
Multitask w/ R+L   & \texttt{<CoT>}   & 85.84 & 75.57 \\
                   & \texttt{<Label>} & \textbf{87.90} & \textbf{77.78} \\ \hline
Multitask w/ C+R+L & \texttt{<CoT>}   & 86.89 & 77.12 \\
                   & \texttt{<Label>} & \textbf{88.23} & \textbf{77.10} \\ \hline
\end{tabular}
\end{table}









\bibliographystyle{IEEEtran}
\bibliography{mybib}

\end{document}

%%% Local Variables:
%%% mode: LaTeX
%%% TeX-master: t
%%% End:
% *==================================================================================*
% *                     Review vs. Camera-Ready settings                             *
% *==================================================================================*
%
% REVIEW: Use the following command for submitting the paper (double-blind,
% for review):
% \documentclass{Interspeech}
%
% CAMERA-READY: Use the following command for the camera-ready version, one
% affiliation per line:
\documentclass[cameraready]{Interspeech}
% *==================================================================================*


% **************************************
% *                                    *
% *      STOP !   DO NOT DELETE !      *
% *          READ THIS FIRST           *
% *                                    *
% * This template also includes        *
% * important INSTRUCTIONS that you    *
% * must follow when preparing your    *
% * paper. Read it BEFORE replacing    *
% * the content with your own work.    *
% **************************************

%==================================================================================
% Title
% Must exactly match the title entered into the paper submission system
\title{Distilling Structured Reasoning into SpeechLLMs for Spoken Language Understanding}

%==================================================================================
% Authors
% The order of authors here must exactly match the order entered into the paper submission system
% Note that the COMPLETE list of authors MUST be entered into the paper submission system at the outset, including when submitting your manuscript for double-blind review
% The ORCID number is still optional but will become mandatory in the future years. It is strongly encouraged to get an ORCID for each cu-author.
% Middle names, including initials, must be included in the first name
\author[affiliation={1}]{Tsukagoshi}{Toshihiro}
\author[affiliation={2}]{Natsuo}{Yamashita}
\author[affiliation={1,3}]{FirstNameC}{LastNameC}
% The maximum number of authors in the author list is 20. If the number of contributing authors is more than this, they should be listed in a footnote or the acknowledgement section.

%==================================================================================
% Affiliations

\address{
    $^1$ Shizuoka University, Japan \\
    $^2$ Address Affiliation 2, Country Affiliation 2 
}

%==================================================================================
% Emails
\email{tsukagoshi.toshihiro.20@shizuoka.ac.jp, second@companyA.com, third@companyB.ai}

%==================================================================================
% Keywords
\keywords{spoken language understanding, speech large language models, chain-of-thought reasoning, multitask learning}

\newcommand{\blue}[1]{\textcolor{blue}{#1}}
\newcommand{\red}[1]{\textcolor{red}{#1}}
\usepackage{comment}

%==================================================================================
% Content

\begin{document}

\maketitle

% the abstract here must exactly match the abstract entered into the paper submission system
\begin{abstract}
Spoken language understanding (SLU) maps speech signals to structured semantic representations, such as intent classification and slot filling. Recent Speech Large Language Models (SpeechLLMs) enable end-to-end SLU by directly processing acoustic inputs. However, standard fine-tuning often exhibits confusion among semantically adjacent intent categories, resulting in insufficient separation between competing interpretations. To mitigate this issue, we propose a structured reasoning distillation framework for SpeechLLM-based SLU. We train the model to learn a stepwise decision process that explicitly contrasts plausible candidate interpretations, rejects incorrect alternatives, and produces the final structured prediction. We generate supervision for this intermediate decision process using a large reasoning model and distill it into SpeechLLMs. Since smaller models may produce unstable structured outputs, we adopt a multitask learning formulation that jointly optimizes intermediate reasoning and final prediction. Experiments on SLURP across six architectures demonstrate consistent improvements over standard fine-tuning.

\end{abstract}






\section{Introduction}

Spoken language understanding (SLU) is a key component of task-oriented dialogue systems, enabling a wide range of real-world applications such as commercial voice assistants~\cite{SaxoChoudharyMcKennaMouchtaris2021E2ESLUVoiceAssistants}, customer-service automation in contact centers~\cite{WuMaoFeng2022AIOnlineCustomerService}, and supporting front-line workers in operational environments~\cite{LinaresGarcia2022VIVA}. 
SLU typically comprises intent classification~\cite{Atuhurra2024DomainAI} and slot filling~\cite{weld2022survey} for mapping an utterance to a semantic frame.


Recent advances in Speech Large Language Models (SpeechLLMs), including SALMONN~\cite{tang2024salmonn}, Qwen2-Audio~\cite{chu2024qwen2audio}, Qwen2.5-Omni~\cite{xu2025qwen25omni}, and Audio-Flamingo~\cite{goel2025audioflamingo3}, enable end-to-end SLU by directly processing acoustic signals and generating structured textual outputs.


Existing studies leveraging SpeechLLMs for SLU have primarily focused on adaptation and generalization.
Some works explore zero-shot or prompting-based SLU by reformulating intent detection and slot filling as question-answering tasks~\cite{LiKeizerDoddipatla2024PromptingWhisperQASLU}, while others evaluate large-scale zero-shot slot filling across diverse schema settings~\cite{hacioglu2025speechllms,hacioglu2025slot}.
Supervised adaptation approaches include parameter-efficient fine-tuning with modality alignment~\cite{li2024whisma}, task-agnostic generalization to unseen intents~\cite{AgrawalGanapathy2025SLUUnseenICL}, and systematic analyses of fine-tuning strategies under limited speech supervision~\cite{choi2025exploring}.
While prior studies have explored the potential of SpeechLLMs for SLU, errors still frequently occur between semantically adjacent intents that differ only in subtle aspects such as scope or granularity.
Learning such fine-grained distinctions purely from data typically requires a substantial number of boundary examples. 
However, collecting sufficiently diverse and finely differentiated training data is costly in practice. Therefore, methods that can efficiently learn decision boundaries from limited data are needed.


An effective alternative to data-intensive boundary learning is to explicitly supervise the intermediate decision process.
In text-based LLM research, rationale-based supervision~\cite{zelikman2022starbootstrappingreasoningreasoning} and distillation methods~\cite{hsieh2023distillingstepbystepoutperforminglarger} have demonstrated that incorporating structured intermediate reasoning can improve downstream performance.
Recent reasoning distillation studies~\cite{magister-etal-2023-teaching, HoSchmidYun2023ReasoningTeachers, Fu2022ComplexityBasedPF} further show that transferring reasoning traces from larger teacher models to smaller students can yield measurable gains under limited model capacity.

However, subsequent analyses suggest that the effectiveness of reasoning supervision depends critically on how it is incorporated.
Naively injecting detailed reasoning traces into smaller models does not always lead to improvements, and excessive reasoning granularity may even hinder optimization~\cite{li-etal-2025-small-models, chen-etal-2025-unveiling-key, chen-etal-2025-skip, hsieh-etal-2023-distilling}.
These findings indicate that reasoning supervision should be treated not as a standalone objective, but as a carefully designed auxiliary signal that shapes the representation space.

Moreover, prior studies have focused primarily on text-only settings, leaving open the question of how structured reasoning supervision behaves in SpeechLLM-based SLU, where acoustic and semantic representations must be jointly modeled.
Motivated by these considerations, we incorporate rationale distillation as an auxiliary objective within a multitask learning framework, allowing structured reasoning to regularize intent prediction without destabilizing generation.



In this work, we investigate whether structured reasoning supervision can improve intent discrimination in moderate-scale SpeechLLMs for end-to-end SLU.
Rather than relying on emergent chain-of-thought behavior, we explicitly model SLU as a stepwise decision process that contrasts plausible intent candidates, rejects incorrect alternatives with evidence-grounded justification, and finally produces a structured semantic label.
We generate such structured rationales using a large reasoning teacher model and distill them into SpeechLLMs under a multitask learning framework, where rationale generation serves as an auxiliary regularizer for direct intent prediction.
Through comprehensive experiments on SLURP across six architectures, we demonstrate that this structured reasoning distillation consistently improves intent accuracy and SLU-F1, while also enlarging inter-class margins in the latent representation space.





\section{Proposed Method}
\label{sec:method}

The overall framework of our approach is illustrated in 
Figure~\ref{fig:framework}. We augment moderate-scale SpeechLLMs 
with structured reasoning to improve performance in low-resource 
SLU settings.

\subsection{Tripartite Rationale Design: $C \to R \to L$}

To transform implicit SLU reasoning into an explicit supervisory 
signal, we define a structured rationale sequence 
$T_{cot} = \{C, R, L\}$. Each component corresponds to a distinct 
reasoning step that mirrors human decision-making under label 
ambiguity.

\begin{itemize}
    \item \textbf{Intent Candidates ($C$)}: Models fine-tuned for SLU 
    often confuse semantically adjacent intents (e.g., 
    \textit{calendar\_set} vs.\ \textit{reminder\_set}), as a 
    direct classification objective provides no explicit incentive 
    to contrast them. To sharpen decision boundaries, we first 
    present a shortlist of plausible intent labels filtered from 
    domain metadata, encouraging contrastive reasoning among the 
    most confusable classes before committing to a prediction.

    \item \textbf{Rejection Rationale ($R$)}: Enumerating candidates 
    alone is insufficient if the model cannot articulate why one is 
    preferred over others. We therefore require a natural language 
    explanation grounded in acoustic or semantic evidence, 
    explicitly justifying the rejection of competing candidates. 
    This step encourages attention to discriminative cues rather 
    than reliance on surface-level lexical patterns.

    \item \textbf{Final Structured Label ($L$)}: After resolving the 
    ambiguity through explicit contrast and justification, the model 
    outputs the final structured semantic representation corresponding 
    to the predicted intent and slot values.
\end{itemize}


By structuring the reasoning in this order, the model is guided 
to progressively narrow the hypothesis space through explicit 
contrastive reasoning, rather than collapsing the decision into 
a single opaque mapping — ultimately producing intent 
representations with greater inter-class separability.
\subsection{Multitask CoT Framework}

A key limitation of moderate-scale SpeechLLMs is that structured 
reasoning does not emerge naturally from scale alone. Rather than 
waiting for such emergence, we explicitly inject reasoning ability 
as a supervised objective through a multitask formulation.

Concretely, the model is conditioned via a mode tag prepended to 
the input, directing it to operate in one of two modes:

\begin{itemize}
    \item \textbf{CoT-mode}: The model generates the full rationale 
    sequence $T_{cot} = \{C, R, L\}$, explicitly performing candidate 
    contrast and rejection before producing the final structured label.

    \item \textbf{Label-mode}: The model directly outputs the final 
    structured label $T_{label} = \{L\}$, without generating intermediate 
    rationale components.
\end{itemize}



Both modes share identical model parameters; the distinction lies 
solely in the generation target specified by the mode tag. Through 
joint training, the reasoning patterns internalized via CoT-mode 
act as a regularizer that shapes the shared representation space, 
enabling Label-mode to produce more discriminative predictions 
without relying on long-form rationale generation at inference time.


\section{Experimental Setup}
\label{sec:experiments}

\subsection{Dataset}
We evaluate our framework on the SLURP 
dataset~\cite{bastianelli-etal-2020-slurp}, which comprises 
11,514 training, 2,033 development, and 2,974 test utterances 
covering 18 domains and 60 intent classes. Following 
\cite{choi2025exploring}, all models are trained on both speech 
and gold transcription inputs to facilitate acquisition of 
fundamental SLU capabilities. For Direct and Standalone CoT 
models, training proceeds for 2 epochs. For Multitask CoT, 
since the model is jointly supervised on both CoT-mode and 
Label-mode targets, the effective number of parameter updates 
per sample is doubled; we therefore train for 1 epoch to ensure 
a fair comparison in terms of total training steps across all 
conditions.

\subsection{Metrics}
Performance is assessed across four complementary metrics. 
Scenario and Action accuracy measure the correctness of 
high-level domain classification (e.g., \textit{calendar}, 
\textit{email}) and specific intent action (e.g., \textit{set}, 
\textit{query}), respectively. Intent accuracy evaluates the 
combined (Scenario, Action) pair, reflecting overall 
classification correctness. SLU-F1~\cite{bastianelli-etal-2020-slurp} 
provides a holistic measure of the full semantic frame, 
penalizing slot errors more heavily when the intent is incorrect.


\begin{enumerate}
    \item \textbf{Scenario}: Accuracy of the high-level domain classification (e.g., \textit{calendar}, \textit{email}).
    \item \textbf{Action}: Accuracy of the specific intent action (e.g., \textit{set}, \textit{query}).
    \item \textbf{Intent}: Accuracy of the combined (Scenario, Action) pair. This represents the overall classification correctness.
    \item \textbf{SLU-F1}: A comprehensive metric that evaluates the entire semantic frame. Unlike standard F1, SLU-F1~\cite{bastianelli-etal-2020-slurp} penalizes slot errors more heavily if the intent is incorrect, providing a holistic measure of the model's understanding of both global intent and local entities.
\end{enumerate}
\begin{table*}[t]
\centering
\caption{Main results on the SLURP test set. Each column group reports 
Scenario (Sce), Action (Act), Intent (Int), and SLU-F1 (F1) accuracy. 
The proposed Multitask CoT consistently achieves the best performance 
across all architectures.}
\label{tab:main_results}
\resizebox{\textwidth}{!}{
\begin{tabular}{lcccccccccccc}
\hline
 & \multicolumn{4}{c}{Direct} 
 & \multicolumn{4}{c}{Standalone CoT (CRL)} 
 & \multicolumn{4}{c}{Multitask CoT (Proposed)} \\ 
\cmidrule(lr){2-5}\cmidrule(lr){6-9}\cmidrule(lr){10-13}
Model & Sce & Act & Int & F1 
      & Sce & Act & Int & F1 
      & Sce & Act & Int & F1 \\ \hline
Qwen2.5-Omni-7B~\cite{xu2025qwen25omni}   & 91.53 & 88.74 & 87.83 & 76.20 & 91.12 & 87.90 & 87.05 & 76.08 & \textbf{91.96} & \textbf{88.97} & \textbf{88.23} & \textbf{77.10} \\
Qwen2.5-Omni-3B~\cite{xu2025qwen25omni}  & 90.65 & 87.79 & 86.79 & 75.04 & 89.27 & 85.47 & 84.47 & 73.23 & \textbf{91.39} & \textbf{88.77} & \textbf{87.79} & \textbf{76.30} \\
Qwen2-Audio-7B~\cite{chu2024qwen2audio}     & 89.37 & 85.21 & 84.23 & 70.08 & 88.63 & 84.80 & 83.46 & 70.62 & \textbf{90.15} & \textbf{86.65} & \textbf{85.71} & \textbf{72.22} \\
Qwen2-Audio-7B-Instruct~\cite{chu2024qwen2audio}   & 88.67 & 85.84 & 84.70 & 70.83 & 88.84 & 85.00 & 83.73 & 71.42 & \textbf{90.11} & \textbf{87.19} & \textbf{86.38} & \textbf{72.79} \\
Audio-Flamingo-3~\cite{goel2025audioflamingo3}    & 87.93 & 84.10 & 82.65 & 69.98 & 87.90 & 84.33 & 83.05 & 71.92 & \textbf{89.88} & \textbf{86.45} & \textbf{85.44} & \textbf{72.86} \\
Music-Flamingo~\cite{ghosh2025musicflamingoscalingmusic}   & 87.69 & 83.89 & 82.65 & 69.02 & 87.96 & 84.36 & 82.92 & 71.42 & \textbf{90.01} & \textbf{86.65} & \textbf{85.74} & \textbf{72.11} \\ \hline
\end{tabular}
}
\end{table*}

\subsection{Representation Analysis}

To quantitatively assess whether multitask CoT supervision improves 
intent discriminability, we extract hidden state vectors at the moment 
of intent generation. Specifically, we identify the tokens corresponding 
to the scenario and action values in the generated output and pool their 
hidden states to obtain a single intent representation vector per 
utterance:

\begin{equation}
    \mathbf{v} = \frac{1}{|\mathcal{T}|} \sum_{t \in \mathcal{T}} \mathbf{h}_t
\end{equation}

\noindent where $\mathcal{T}$ denotes the set of token positions 
corresponding to the scenario and action values, and $\mathbf{h}_t$ 
is the hidden state at position $t$ from the final layer.

We evaluate inter-class separability using two complementary metrics. 
The \textbf{Silhouette coefficient} for sample $i$ is defined as:

\begin{equation}
    s(i) = \frac{b(i) - a(i)}{\max\{a(i),\, b(i)\}}
\end{equation}

\noindent where $a(i)$ is the mean cosine distance to samples in the 
same intent class, and $b(i)$ is the mean cosine distance to samples 
in the nearest competing class. The overall score $\bar{s} = 
\frac{1}{N}\sum_i s(i)$ approaches $1$ when clusters are 
well-separated.

The \textbf{nearest-centroid margin} for sample $i$ is defined as:

\begin{equation}
    m(i) = \min_{c \neq y_i} \lVert\mathbf{v}_i - \boldsymbol{\mu}_c\rVert_2 
           - \lVert\mathbf{v}_i - \boldsymbol{\mu}_{y_i}\rVert_2
\end{equation}

\noindent where $\boldsymbol{\mu}_c$ is the centroid of intent class 
$c$ and $y_i$ is the ground-truth intent of sample $i$. A larger 
positive margin indicates a clearer decision boundary between 
semantically adjacent intents.



\subsection{Implementation Details}
All models are fine-tuned using BF16 precision with the AdamW 
optimizer, a learning rate of $4 \times 10^{-5}$, a batch size 
of 4, and 4 gradient accumulation steps. Oracle rationales for 
multitask training are generated via DeepSeek-R1, whose 
chain-of-thought pretraining makes it well-suited for producing 
structured, evidence-grounded reasoning traces in the 
$C \to R \to L$ format. For Direct and Standalone CoT models, 
training proceeds for 2 epochs. For Multitask CoT, since the 
model is jointly supervised on both CoT-mode and Label-mode 
targets, the effective number of parameter updates per sample 
is doubled; we therefore train for 1 epoch to ensure a fair 
comparison in terms of total training steps across all conditions.

\begin{table}[t]
\centering
\caption{Representation geometry metrics for Qwen2.5-Omni-7B. 
Higher values indicate greater inter-class separability.}
\label{tab:rep_analysis}
\resizebox{\columnwidth}{!}{
\begin{tabular}{lccc}
\hline
\textbf{Metric} & \textbf{Direct} & \textbf{Multitask CoT} & \textbf{Diff} \\ \hline
Silhouette (cosine)          & 0.242 & 0.279 & +0.037 \\
Silhouette (euclidean)       & 0.157 & 0.176 & +0.019 \\
Fisher ratio                 & 1.133 & 1.276 & +0.143 \\
Inter-centroid cosine dist.  & 0.262 & 0.328 & +0.067 \\
Norm. centroid margin (mean) & 0.174 & 0.194 & +0.020 \\ \hline
\end{tabular}
}
\end{table}

\section{Results}
\label{sec:results}

\subsection{Main Results}
Table~\ref{tab:main_results} reports Scenario, Action, Intent accuracy and SLU-F1 for three training conditions across six architectures: Direct fine-tuning, Standalone CoT, and the proposed Multitask CoT using \texttt{<Label>} mode inference.
Multitask CoT consistently outperforms both baselines across all models and metrics, with intent accuracy gains ranging from +0.96\% (Qwen2.5-Omni-7B) to +2.79\% (Qwen2.5-Omni-3B) over Direct fine-tuning, and SLU-F1 improvements of up to +1.26 points (Qwen2.5-Omni-7B).
Notably, Standalone CoT occasionally underperforms Direct fine-tuning, with intent accuracy degrading by up to 2.32\% (Qwen2.5-Omni-3B), indicating that CoT supervision without a complementary \texttt{<Label>} mode objective does not reliably improve performance in moderate-scale models.
These results demonstrate that explicitly supervising reasoning as a multitask objective is more effective than relying on emergent CoT behavior for moderate-scale SpeechLLMs.


\subsection{Representation Analysis}

Table~\ref{tab:rep_analysis} reports representation geometry 
metrics comparing Direct fine-tuning and Multitask CoT 
(Qwen2.5-Omni-7B). Multitask CoT improves across all 
separability metrics: the Silhouette coefficient (cosine) 
increases from 0.242 to 0.279 (+0.037), and the Fisher ratio 
rises from 1.133 to 1.276 (+0.143), indicating greater 
between-class scatter relative to within-class scatter. The 
mean inter-centroid cosine distance also increases from 0.262 
to 0.328 (+0.067), suggesting that intent clusters are more 
evenly spread across the representation space. Furthermore, 
the normalized nearest-centroid margin improves from 0.174 to 
0.194 (+0.020), confirming that decision boundaries between 
semantically adjacent intents become more clearly defined 
under multitask CoT supervision.



\subsection{Impact of Rationale Components}
Table~\ref{tab:ablation_components} reports intent accuracy and SLU-F1 for each rationale component configuration using \texttt{<Label>} mode inference. Among individual components, C yields the largest improvement in intent accuracy (+0.37\% over baseline), indicating that explicitly enumerating candidate intents contributes most to sharpening decision boundaries. R alone provides marginal gains, suggesting that rejection rationale is insufficient as a standalone supervisory signal. The full configuration C+R+L achieves the highest intent accuracy, while C+L and C+R+L both attain the highest SLU-F1 of 77.18, indicating that the combination of candidate enumeration and rejection rationale contributes to stable performance across metrics.



\begin{table}[t]
\centering
\caption{Ablation of rationale components (Qwen2.5-Omni-7B, Label-branch inference). \textbf{C}: Intent Candidates, \textbf{R}: Rejection Rationale, \textbf{L}: Label (direct prediction).}
\label{tab:ablation_components}
\begin{tabular}{lcccc}
\hline
\textbf{Configuration} & \textbf{Sce / Act / Int} & \textbf{SLU-F1} \\ \hline
Direct (Baseline)      & 91.53 / 88.74 / 87.83 & 76.20 \\
Multitask w/ C         & 91.56 / 89.00 / 88.20 & 77.18 \\
Multitask w/ R         & 91.53 / 88.84 / 87.96 & 76.01 \\
Multitask w/ C+R       & 91.83 / 88.87 / 88.13 & \textbf{77.18} \\
\textbf{Multitask w/ C+R+L} & \textbf{91.96 / 88.97 / 88.23} & 77.10 \\ \hline
\end{tabular}
\end{table}

The results indicate that providing Intent Candidates ($C$) yields the most significant initial boost in intent accuracy (+0.37\% over baseline). This supports our hypothesis that forcing the model to explicitly contrast similar intent categories early in the process sharpens the decision boundaries. While Rejection Rationale ($R$) alone provides marginal gains, its combination with candidates (CRL) achieves the highest overall accuracy and stability across domains.

\subsection{Inference Branch: CoT vs. Label}
Table~\ref{tab:branch_comparison} reports intent accuracy and SLU-F1 for each multitask configuration under two inference modes: \texttt{<CoT>} and \texttt{<Label>}. Across all configurations, \texttt{<Label>} mode consistently achieves higher performance than \texttt{<CoT>} mode. Notably, the performance gap is particularly pronounced in the Multitask w/ R+L configuration, where \texttt{<Label>} mode outperforms \texttt{<CoT>} mode by 2.06\% in intent accuracy. This trend suggests that CoT-mode supervision effectively shapes the shared representation space during training, but the benefit does not transfer to inference-time rationale generation. Rather, the structured reasoning internalized through CoT-mode training is more effectively utilized when the model performs direct prediction via \texttt{<Label>} mode, without the overhead of generating intermediate rationale sequences.


\section{Conclusion}
\label{sec:conclusion}
In this paper, we proposed a Multitask CoT framework for 
moderate-scale SpeechLLMs to address the challenges of spoken 
language understanding. By supervising a structured reasoning process 
($C \to R \to L$) alongside direct intent prediction, we successfully 
improved performance across six distinct model architectures. 
Our representation analysis, including Silhouette coefficients and 
Fisher ratio, provided quantitative evidence that multitask rationale 
supervision encourages the model to develop more separable intent 
clusters. Crucially, we demonstrated that using the rationale as a 
training regularizer while performing direct inference (\texttt{<Label>} 
mode) allows the model to leverage the benefits of reasoning without 
the overhead or noise of long-form generation. Future work will explore 
the scalability of this approach to even smaller edge-device models and 
its application to multilingual SLU scenarios.

\begin{table}[t]
\centering
\caption{Comparison of inference modes within the multitask framework.}
\label{tab:branch_comparison}
\begin{tabular}{llcc}
\hline
\textbf{Configuration} & \textbf{Mode} & \textbf{Intent Acc.} & \textbf{SLU-F1} \\ \hline
Multitask w/ C+L   & \texttt{<CoT>}   & 87.66 & 77.31 \\
                   & \texttt{<Label>} & \textbf{88.06} & \textbf{77.89} \\ \hline
Multitask w/ R+L   & \texttt{<CoT>}   & 85.84 & 75.57 \\
                   & \texttt{<Label>} & \textbf{87.90} & \textbf{77.78} \\ \hline
Multitask w/ C+R+L & \texttt{<CoT>}   & 86.89 & 77.12 \\
                   & \texttt{<Label>} & \textbf{88.23} & \textbf{77.10} \\ \hline
\end{tabular}
\end{table}









\bibliographystyle{IEEEtran}
\bibliography{mybib}

\end{document}

%%% Local Variables:
%%% mode: LaTeX
%%% TeX-master: t
%%% End:
